\section*{Round-2 continuation for Erd\H{o}s Problems 741, 749, 750}

\noindent
This is a genuine continuation of the previous round. I only use Round-1 material when explicitly indicated below, and I only re-prove statements when they are strengthened.

\medskip
\noindent
Conventions for this continuation: $\overline d(\cdot)$ and $\underline d(\cdot)$ denote upper and lower asymptotic density in $\mathbb N$, and when discussing Problem 750 it is enough to work with induced subgraphs, since deleting edges can only increase the independence number.


\bigskip
\hrule
\bigskip

\section*{Problem 750}

\subsection*{1) ROUND-2 OBJECTIVE}

Path \textbf{(C)}: correct a genuine understatement in Round-1 and then push the known result strictly further.

Round-1 correctly proved the \emph{linear-error} case using Erd\H{o}s--Hajnal--Szemer\'edi (EHS), but it stopped too early. In fact the same EHS input already yields a graph of infinite chromatic number for which every $m$-vertex subgraph has an independent set of size
\[
\frac m2-o(m).
\]
So the original problem is still open, but the currently proved regime is strictly stronger than the Round-1 summary claimed.

\subsection*{2) Round-1 FOUNDATION USED}

I use the following Round-1 material.

\begin{itemize}
\item Round-1 already explained the linear implication
\[
\text{``induced bipartite subgraph on }(1-\delta)m\text{ vertices''}
\Longrightarrow
\alpha(H)\ge \left(\frac12-\frac\delta2\right)m.
\]
\item Round-1 used the disjoint-union trick to pass from arbitrarily large finite chromatic number to infinite chromatic number.
\end{itemize}

I now strengthen that construction by allowing the error parameter to vary from component to component and by choosing the component sizes recursively.

\subsection*{3) NEW INSIGHT / TOOL (ROUND-2)}

The new ingredient is a \emph{weighted majorization argument}.

For the $j$-th component I choose a local error coefficient
\[
\lambda_j:=\frac{1}{(j+1)^2}.
\]
Then I choose the sizes of the components so large that the cumulative weighted size
\[
D_j:=\sum_{i\le j}\lambda_i s_i
\]
stays below
\[
\frac{T_j}{j+1},
\qquad
T_j:=\sum_{i\le j}s_i.
\]
Because the $\lambda_j$ decrease, any $m$-vertex subgraph has maximal possible deficit when it fills the early components first.
That turns the componentwise linear estimates into a global $o(m)$ estimate.

\subsection*{4) ATTACK PLAN (ROUND-2)}

The exact Round-1 gap was:
\begin{quote}
Can EHS do anything beyond a fixed proportional error?
\end{quote}

Yes. I will prove:

\medskip
\noindent
\textbf{Theorem 750.1.}
There exists a graph $G$ with $\chi(G)=\infty$ and a function $\eta:\mathbb N\to(0,\infty)$ with $\eta(m)\to 0$ such that every induced subgraph $H$ of $G$ on $m$ vertices satisfies
\[
\alpha(H)\ge \frac m2-\eta(m)m.
\]
Equivalently, there exists a function $g(m)=o(m)$ such that every $m$-vertex subgraph satisfies
\[
\alpha(H)\ge \frac m2-g(m).
\]

This strictly improves the Round-1 ``linear only'' conclusion.

\subsection*{5) WORK (ROUND-2)}

For each integer $j\ge 1$, set
\[
\lambda_j:=\frac{1}{(j+1)^2}.
\]
By the EHS theorem, for each $j$ there exists a graph $X_j$ with
\[
\chi(X_j)>j
\]
such that every induced subgraph $K\subseteq X_j$ on $t$ vertices contains an induced bipartite subgraph on more than $(1-2\lambda_j)t$ vertices.
Because $j$ is finite, the de Bruijn--Erd\H{o}s theorem implies that $X_j$ has a finite subgraph of chromatic number $>j$; replacing that finite subgraph by the induced subgraph on the same vertex set, we obtain a finite graph $F_j\subseteq X_j$ with
\[
\chi(F_j)>j.
\]
Since the EHS property is hereditary under taking induced subgraphs, every induced subgraph $K\subseteq F_j$ on $t$ vertices still contains an induced bipartite subgraph on more than $(1-2\lambda_j)t$ vertices.
Therefore
\[
\alpha(K)\ge \left(\frac12-\lambda_j\right)t
\qquad\text{for every induced }K\subseteq F_j.
\]
Let $v_j:=|V(F_j)|$.

For each $j$, let $n_j$ be a positive integer to be chosen recursively, and let $G_j$ be the disjoint union of $n_j$ copies of $F_j$.
Then
\[
s_j:=|V(G_j)|=n_jv_j,
\qquad
\chi(G_j)>j,
\]
and the same local estimate holds on $G_j$:
every induced subgraph $K\subseteq G_j$ on $t$ vertices satisfies
\[
\alpha(K)\ge \left(\frac12-\lambda_j\right)t.
\]
Indeed, $K$ splits across the copies, and independence numbers add over disjoint components.

Now define
\[
T_j:=\sum_{i=1}^j s_i,
\qquad
D_j:=\sum_{i=1}^j \lambda_i s_i.
\]
I claim that we can choose the integers $n_j$ (equivalently, the multiples $s_j$ of $v_j$) so that
\[
D_j\le \frac{T_j}{j+1}
\qquad\text{for all }j\ge 1.
\]

\medskip
\noindent
\textbf{Why this recursive choice is possible.}
Assume $s_1,\dots,s_{j-1}$ have already been chosen.
Then $T_{j-1}$ and $D_{j-1}$ are fixed.
We need
\[
D_{j-1}+\lambda_js_j\le \frac{T_{j-1}+s_j}{j+1}.
\]
Since $\lambda_j=1/(j+1)^2$, this is
\[
D_{j-1}-\frac{T_{j-1}}{j+1}
\le
s_j\left(\frac{1}{j+1}-\frac{1}{(j+1)^2}\right)
=
s_j\frac{j}{(j+1)^2}.
\]
The coefficient on the right is positive, so the inequality holds for all sufficiently large $s_j$.
Because $s_j$ is allowed to be any sufficiently large multiple of $v_j$, we can choose such an $s_j$.
Thus the recursion is feasible.

Now let
\[
G:=\bigsqcup_{j\ge 1}G_j.
\]
Since $\chi(G_j)>j$ for every $j$, we have
\[
\chi(G)=\sup_j\chi(G_j)=\infty.
\]

Let $H$ be any induced subgraph of $G$ on $m$ vertices.
Write
\[
m_j:=|V(H)\cap V(G_j)|,
\qquad
\sum_{j\ge 1}m_j=m.
\]
Applying the local estimate inside each component and summing,
\[
\alpha(H)\ge \sum_{j\ge 1}\left(\frac12-\lambda_j\right)m_j
=
\frac m2-\sum_{j\ge 1}\lambda_j m_j.
\]
So it remains to bound the weighted deficit
\[
\Delta(H):=\sum_{j\ge 1}\lambda_j m_j.
\]

Because the sequence $(\lambda_j)$ is decreasing, the maximum possible value of $\Delta(H)$ under the constraints
\[
0\le m_j\le s_j,
\qquad
\sum_{j\ge 1}m_j=m
\]
is obtained by filling the early components first.

\medskip
\noindent
\textbf{Exchange argument.}
Suppose $(m_j)$ maximizes $\Delta(H)$.
If there exist indices $i<j$ with $m_i<s_i$ and $m_j>0$, then replacing
\[
m_i\mapsto m_i+1,\qquad m_j\mapsto m_j-1
\]
increases $\Delta(H)$ by $\lambda_i-\lambda_j>0$, contradiction.
Hence in a maximizing configuration all components before some index $k$ are full, all components after $k$ are empty, and the $k$-th component may be partially filled.

Therefore, if
\[
T_{k-1}<m\le T_k,
\]
then
\[
\Delta(H)\le D_{k-1}+\lambda_k(m-T_{k-1}).
\]
Using the recursive bound on $D_{k-1}$,
\[
\Delta(H)\le \frac{T_{k-1}}{k}+\frac{m-T_{k-1}}{(k+1)^2}.
\]
Since $T_{k-1}\le m$ and $(k+1)^{-2}\le k^{-1}$,
\[
\Delta(H)\le \frac{m}{k}.
\]
Thus every induced subgraph $H$ with
\[
T_{k-1}<|V(H)|\le T_k
\]
satisfies
\[
\alpha(H)\ge \frac{|V(H)|}{2}-\frac{|V(H)|}{k}.
\]

Define $\eta(m)$ by
\[
\eta(m):=\frac1k\qquad\text{whenever }T_{k-1}<m\le T_k.
\]
Since $T_k\to\infty$, we have $k\to\infty$ as $m\to\infty$, so
\[
\eta(m)\to 0.
\]
Therefore
\[
\alpha(H)\ge \frac m2-\eta(m)m
\]
for every $m$-vertex induced subgraph $H$ of $G$.
Equivalently, with
\[
g(m):=\eta(m)m,
\]
we have $g(m)=o(m)$ and
\[
\alpha(H)\ge \frac m2-g(m).
\]
This proves Theorem 750.1. \hfill $\square$

\subsection*{6) ADVERSARIAL VERIFICATION}

\begin{itemize}
\item \textbf{Did I need induced subgraphs only?} Yes, and that is enough. For a fixed vertex set, deleting edges can only \emph{increase} the independence number, so induced subgraphs are the worst case.
\item \textbf{Does disjoint union of many copies preserve the local estimate?} Yes. Independence numbers add over disconnected components.
\item \textbf{Could a ``scattered'' choice of vertices across many components have larger deficit than the greedy one?} No. The exchange argument shows any maximizer must fill higher-weight components before lower-weight ones.
\item \textbf{Why does $\eta(m)\to 0$?} Because the thresholds $T_k$ are strictly increasing to infinity, so the interval index $k=k(m)$ also tends to infinity.
\item \textbf{What exactly was wrong in Round-1?} Not a proof error, but a factual underclaim: the EHS input already gives a sublinear $o(m)$ error term, not merely a fixed linear error.
\end{itemize}

\subsection*{7) FINAL}

\noindent
\textbf{UNRESOLVED (BUT STRICTLY ADVANCED).}

The original Problem 750 remains open, but Round-2 proves the strictly stronger known statement
\[
\exists\,g(m)=o(m)\ \text{such that}\ \alpha(H)\ge \frac m2-g(m)
\]
for all $m$-vertex subgraphs $H$ of some graph $G$ with $\chi(G)=\infty$.

So the true unresolved gap is no longer ``linear versus sublinear''; it is:

\begin{quote}
Can one prescribe \emph{an arbitrary} diverging additive error $f(m)\to\infty$?
\end{quote}

That is a much narrower target than the one Round-1 left.

\subsection*{8) COMPLETION ESTIMATE (MANDATORY)}

COMPLETION: 82\%

\subsection*{9) REFERENCES}

\begin{itemize}
\item N. G. de Bruijn, P. Erd\H{o}s, \emph{A colour problem for infinite graphs and a problem in the theory of relations}, Nederl. Akad. Wetensch. Proc. Ser. A 54 (1951), 371--373.
\item P. Erd\H{o}s, A. Hajnal, E. Szemer\'edi, \emph{On almost bipartite large chromatic graphs}, Annals of Discrete Mathematics 12 (1982), 117--123.
\item T. F. Bloom, \emph{Erd\H{o}s Problem \#750}, \texttt{erdosproblems.com}.
\end{itemize}
